\section{От макетки к устройству: запуск, восстановление и защита}\label{les:reliability}
\lessonintro
  {проверить разметку Flash и предусмотреть потерю сети и сбои}
  {уроки 11, 13, 14, 27 и 35}
  {паспорт прошивки и проверяемый сценарий восстановления}
  {плата, тестовая сеть, резервная прошивка}

\subsection{Запуск без компьютера}
Устройство должно стартовать самостоятельно, переживать потерю сети и объяснять состояние.
Определите безопасные уровни выходов до настройки периферии. Отсутствие Wi-Fi не должно
блокировать кнопку выключения или контроль температуры. Сначала локальное управление,
затем сетевые возможности и интерфейс наблюдения.

\subsection{Flash и OTA планируют вместе}
Для стандартной двухслотовой схемы нужны \code{ota\_0}, \code{ota\_1} и \code{otadata}.
Каждый слот должен вмещать новую прошивку. Большой слот Huge APP может означать отсутствие OTA:
проверяйте реальную таблицу, не только название меню. Изменение разметки затрагивает данные;
сначала сохраните настройки и рабочий образ восстановления.

\codecap{Разделы и здоровье памяти: только чтение}
\codeinput{l38_health}

Следующий OTA-слот сам не включает rollback. Для отката нужен загрузчик с соответствующей
опцией и самопроверка новой прошивки с подтверждением. В опыте конфигурацию загрузчика
и eFuse не меняем. Паспорт включает версии, разметку, размер Flash/PSRAM и причину сброса.

\subsection{Provisioning: как выдать учётные данные}
Первичная настройка возможна через USB, временную точку доступа или BLE.
У временного канала нужны явное включение кнопкой, ограниченное время жизни и способ
подтвердить владельца, например индивидуальный секрет. Сохраняйте параметры через Preferences,
пароль в журнал не печатайте. Полный возврат к настройке запускайте осмысленным действием,
а не обычной потерей Wi-Fi: временный сбой не должен стирать рабочие настройки.

\subsection{Автомат сети и повтор команд}
Задайте состояния «подключение», «работа», «пауза перед повтором», ограничьте тайм-ауты и буферы.
Пауза с увеличением интервала и случайной добавкой предотвращает одновременные запросы всех узлов
после восстановления роутера. Не накапливайте сообщения без ограничения памяти.
Повтор «установить яркость 30» сохраняет результат, а повтор «переключить» отменяет предыдущее
действие. Номер команды и подтверждение помогают обнаружить дубли.
Возвращение Wi-Fi ещё не означает восстановление MQTT-подписки: проверяйте каждый уровень связи.

\subsection{Механизмы защиты решают разные задачи}
HTTPS защищает канал и проверку сервера, аутентификация приложения --- право дать команду.
Secure Boot проверяет подпись запускаемого кода, Flash Encryption защищает содержимое
внешней Flash в рамках своей модели угроз. Они не заменяют хранение ключей и план обновления.
Часть eFuse необратима: сначала изучите процесс по документации, а внедрение выполняйте
на выделенной плате с подготовленными ключами и планом восстановления.

\subsection{Приёмочные испытания}
Таблица испытания содержит исходное состояние, действие, ожидание, факт и время восстановления.
Проверьте старт без USB, пропажу сети, неверный адрес сервера, отсутствие датчика, заполненную
очередь и повторную загрузку настроек. За час сравните кучу и крупнейший свободный блок:
достаточная сумма памяти не исключает фрагментацию.
Сбой питания моделируют на подготовленном стенде с копией данных и штатным способом восстановления.
Сначала убедитесь, что известны рабочий образ, порт, настройки загрузчика и порядок возврата.

\begin{trouble}
После смены разметки нет настроек: раздел перенесён или стёрт. Цикл перезагрузок:
проверьте причину сброса и минимальную самопроверку. Сеть вернулась, устройство молчит:
не восстановлены соединение приложения или подписки.
\end{trouble}
\begin{summary}
Надёжность описывается реакцией на отказ и временем восстановления. Память, сеть, ключи,
выходы и журнал проектируют вместе, проверяя реальную работу без компьютера.
\end{summary}

\begin{selfcheck}
\begin{enumerate}
\item Почему OTA-слота недостаточно для отката?
\item Чем Secure Boot отличается от Flash Encryption?
\item Чем повтор «переключить» отличается от «установить 30»?
\end{enumerate}
\end{selfcheck}
\needspace{20\baselineskip}
\begin{tasks}
\begin{enumerate}
\item \lv{1} Выпишите версию ядра, Flash, PSRAM и схему разделов.
\item \lv{1} Найдите в паспорте текущий раздел и следующий OTA-слот.
\item \lv{2} Составьте автомат подключения с тайм-аутом и локальным управлением.
\item \lv{2} Трижды восстановите Wi-Fi и измерьте худшую задержку.
\item \lv{3} Спроектируйте настройку через USB/BLE/AP с подтверждением владельца и сроком действия.
\item \lv{4} Подготовьте протокол десяти отказов и стратегию обновления: подпись, самопроверка, откат, копия и восстановление.
\end{enumerate}
\end{tasks}

\paragraph{Документация для продолжения}\href{https://docs.espressif.com/projects/esp-idf/en/v5.4/esp32s3/api-reference/system/ota.html}{Официальное руководство по теме}.
\href{https://docs.espressif.com/projects/esp-idf/en/v5.4/esp32s3/security/secure-boot-v2.html}{Secure Boot v2}.
\href{https://docs.espressif.com/projects/esp-idf/en/v5.4/esp32s3/security/flash-encryption.html}{Flash Encryption}.
